%%%%%%%%%%%%%%%%
%%% deut 22 %%%
%%%%%%%%%%%%%%%%

\Chapter{22}

\Verse{1}
{
	\hDtn{לֹא תִרְאֶה אֶת שׁוֹר אָחִיךָ אוֹ אֶת שֵׂיוֹ נִדָּחִים וְהִתְעַלַּמְתָּ מֵהֶם הָשֵׁב תְּשִׁיבֵם לְאָחִיךָ׃}
}
{
	\eDtn{
		“You shall not see your brother’s ox or his sheep
		driven off, and you hide yourself from them. You shall
		bring them back to your brother.
	}
}
{
}

\Verse{2}
{
	\hDtn{וְאִם לֹא קָרוֹב אָחִיךָ אֵלֶיךָ וְלֹא יְדַעְתּוֹ וַאֲסַפְתּוֹ אֶל תּוֹךְ בֵּיתֶךָ וְהָיָה עִמְּךָ עַד דְּרֹשׁ אָחִיךָ אֹתוֹ וַהֲשֵׁבֹתוֹ לוֹ׃}
}
{
	\eDtn{
		And if your brother is not close to you, and you don’t
		know him, then you shall gather it into your house, and it
		shall be with you until your brother inquires about it,
		and you shall give it back to him.
	}
}
{
}

\Verse{3}
{
	\hDtn{וְכֵן תַּעֲשֶׂה לַחֲמֹרוֹ וְכֵן תַּעֲשֶׂה לְשִׂמְלָתוֹ וְכֵן תַּעֲשֶׂה לְכל אֲבֵדַת אָחִיךָ אֲשֶׁר תֹּאבַד מִמֶּנּוּ וּמְצָאתָהּ לֹא תוּכַל לְהִתְעַלֵּם׃}
}
{
	\eDtn{
		And you shall do that with his ass, and you shall do that
		with his garment, and you shall do that with any lost
		thing of your brother’s that will be lost by him and you
		find it. You may not hide yourself.
	}
}
{
}

\Verse{4}
{
	\hDtn{לֹא תִרְאֶה אֶת חֲמוֹר אָחִיךָ אוֹ שׁוֹרוֹ נֹפְלִים בַּדֶּרֶךְ וְהִתְעַלַּמְתָּ מֵהֶם הָקֵם תָּקִים עִמּוֹ׃}
}
{
	\eDtn{
		You shall not see your brother’s ass or his ox falling in
		the road, and you hide yourself from them. You shall lift
		it up with him.
	}
}
{
}

\Verse{5}
{
	\hDtn{לֹא יִהְיֶה כְלִי גֶבֶר עַל אִשָּׁה וְלֹא יִלְבַּשׁ גֶּבֶר שִׂמְלַת אִשָּׁה כִּי תוֹעֲבַת יְהֹוָה אֱלֹהֶיךָ כּל עֹשֵׂה אֵלֶּה׃}
}
{
	\eDtn{
		“There shall not be a man’s item on a woman, and a man
		shall not wear a woman’s garment, because everyone who
		does these is an offensive thing of YHWH, your God.
	}
}
{
}

\Verse{6}
{
	\hDtn{כִּי יִקָּרֵא קַן צִפּוֹר לְפָנֶיךָ בַּדֶּרֶךְ בְּכל עֵץ אוֹ עַל הָאָרֶץ אֶפְרֹחִים אוֹ בֵיצִים וְהָאֵם רֹבֶצֶת עַל הָאֶפְרֹחִים אוֹ עַל הַבֵּיצִים לֹא תִקַּח הָאֵם עַל הַבָּנִים׃}
}
{
	\eDtn{
		“When a bird’s nest will happen to be in front of you on
		the road in any tree or on the ground—chicks or eggs—and
		the mother is sitting over the chicks or over the eggs,
		you shall not take the mother along with the children.
	}
}
{
%	\footnote{
%		the mother along with the children. The purpose of this law, like the
%		law against cooking a kid in its mother’s milk, is uncertain. Both laws
%		may have to do with caring about animals, prohibiting unfeeling
%		treatment of them. But both may also have purposes of a different sort.
%		Both cases may have been regarded as unnatural combinations, offensive
%		in themselves. In the case of the mother bird along with the eggs or
%		chicks, that would explain why this law is grouped here with laws
%		against forbidden mixtures of seeds, fabrics, and work animals, as well
%		as the law against males wearing female clothing and females wearing
%		male items. (For another explanation of the kid cooked in milk, see the
%		comment on 22:9 below.)
%	}
}

\Verse{7}
{
	\hDtn{שַׁלֵּחַ תְּשַׁלַּח אֶת הָאֵם וְאֶת הַבָּנִים תִּקַּח לָךְ לְמַעַן יִיטַב לָךְ וְהַאֲרַכְתָּ יָמִים׃}
}
{
	\eDtn{
		You shall let the mother go, and you may take the children
		for you, so that it will be good for you, and you will
		extend days.
	}
}
{
}

\Verse{8}
{
	\hDtn{כִּי תִבְנֶה בַּיִת חָדָשׁ וְעָשִׂיתָ מַעֲקֶה לְגַגֶּךָ וְלֹא תָשִׂים דָּמִים בְּבֵיתֶךָ כִּי יִפֹּל הַנֹּפֵל מִמֶּנּוּ׃}
}
{
	\eDtn{
		“When you’ll build a new house, you shall make a railing
		for your roof, so you won’t set blood in your house when
		someone will fall from it.
	}
}
{
%	\footnote{
%		a railing for your roof. This law seems to be out of place, in the
%		middle of laws concerning mixtures. But there is apparently more than
%		one principle of orproduce will beganization operating here. The first
%		laws of this chapter concern caring about one’s neighbor, but
%		specifically expressed in terms of the neighbor’s animals. Then it adds
%		the neighbor’s garment. These laws are followed by the law prohibiting
%		men and women to mix each other’s garments, followed by a return to laws
%		regarding animals, followed by another law of caring about one’s
%		neighbor: this law of putting a railing on one’s roof. This is followed
%		by more laws of prohibited mixtures, including mixing certain animals.
%		The full group is a complex mixture of moral and ritual laws, but as a
%		group it also shows how the Torah’s different kinds of laws are
%		interconnected. Footnote 22:8. set blood in your house. Meaning: you
%		won’t bring the responsibility on your home for a loss of life.
%	}
}

\Verse{9}
{
	\hDtn{לֹא תִזְרַע כַּרְמְךָ כִּלְאָיִם פֶּן תִּקְדַּשׁ הַמְלֵאָה הַזֶּרַע אֲשֶׁר תִּזְרָע וּתְבוּאַת הַכָּרֶם׃}
}
{
	\eDtn{
		“You shall not seed your vineyard with two kinds, or else
		the whole of the seed that you’ll sow and the vineyard’s
		produce will become holy.
	}
}
{
%	\footnote{
%		will become holy. Why does mixing different seeds in a vineyard make the
%		resulting crop holy? Several of the laws regarding mixtures have
%		possible explanations having to do with the realm of the holy. The law
%		against cooking a kid in its mother’s milk may be because that was
%		regarded as a food for a deity, since a Ugaritic text pictures the chief
%		god, El, having kid cooked in milk (see the comment on Exod 23:19). The
%		law against wearing wool and linen together may be because they were
%		both used in the Tabernacle (see the comment on Deut 22:11 below). (Even
%		the law against men wearing clothes of the opposite sex may have to do
%		with pagan myths in which the gods and goddesses do that.) And so it may
%		be in the case of mixed seeds, as well: the prohibition of mixing them
%		may not be because the mixing is bad in some way but rather because some
%		mixtures are regarded as divine. This is only a speculation, since the
%		meaning of this law (and of the other laws concerning mixtures) has
%		stymied scholars for centuries. But this speculation at least has the
%		advantage of coming to terms with the word “holy” in this law. Other
%		commentaries and translations have commonly ascribed some other meaning
%		to the word because of the difficulty of making sense of it in this
%		context.
%	}
}

\Verse{10}
{
	\hDtn{לֹא תַחֲרֹשׁ בְּשׁוֹר וּבַחֲמֹר יַחְדָּו׃}
}
{
	\eDtn{You shall not plow with an ox and an ass together.}
}
{
}

\Verse{11}
{
	\hDtn{לֹא תִלְבַּשׁ שַׁעַטְנֵז צֶמֶר וּפִשְׁתִּים יַחְדָּו׃}
}
{
	\eDtn{You shall not wear sha‘atnez: wool and linen together.}
}
{
%	\footnote{
%		sha‘atnez: wool and linen together. The former comes from an animal, a
%		sheep; and the latter comes from a plant, flax. Some think that such
%		fabric is prohibited because it was thought to be an unnatural mixture.
%		Others think that it is because the priests have both linen and wool in
%		their clothing, and that therefore laypersons must not wear what belongs
%		to the realm of the sacred. The problem with the latter view is that the
%		only description of priestly garments in the Torah mentions linen but
%		does not mention wool explicitly. Still, it is likely that the priests
%		have wool and linen, because some of their clothing is said to be dyed
%		(Exod 28:4–6), and it is extremely difficult to dye linen (I learned
%		this from Avigail Sheffer, a specialist in ancient textiles). Fabric
%		excavated at Kuntillat ‘Ajrud, which may have been a cultic site,
%		contained linen and wool. Alternatively, the Tabernacle is made of an
%		inside layer of linen fabric and a second layer of wool (goats’ hair)
%		fabric over it. The sha‘atnez prohibition may therefore relate to the
%		sacred state of the Tabernacle rather than of the priesthood.
%	}
}

\Verse{12}
{
	\hDtn{גְּדִלִים תַּעֲשֶׂה לָּךְ עַל אַרְבַּע כַּנְפוֹת כְּסוּתְךָ אֲשֶׁר תְּכַסֶּה בָּהּ׃}
}
{
	\eDtn{
		“You shall make braided threads on the four corners of
		your apparel with which you cover yourself.
	}
}
{
}

\Verse{13}
{
	\hDtn{כִּי יִקַּח אִישׁ אִשָּׁה וּבָא אֵלֶיהָ וּשְׂנֵאָהּ׃}
}
{
	\eDtn{
		“When a man will take a wife and come to her and then hate
		her,
	}
}
{
}

\Verse{14}
{
	\hDtn{וְשָׂם לָהּ עֲלִילֹת דְּבָרִים וְהוֹצִא עָלֶיהָ שֵׁם רָע וְאָמַר אֶת הָאִשָּׁה הַזֹּאת לָקַחְתִּי וָאֶקְרַב אֵלֶיהָ וְלֹא מָצָאתִי לָהּ בְּתוּלִים׃}
}
{
	\eDtn{
		and he’ll assert words of abuse toward her and bring out a
		bad name on her and say, ‘I took this woman, and I came
		close to her, and I didn’t find signs of virginity for
		her,’
	}
}
{
}

\Verse{15}
{
	\hDtn{וְלָקַח אֲבִי הַנַּעֲרָ וְאִמָּהּ וְהוֹצִיאוּ אֶת בְּתוּלֵי הַנַּעֲרָ אֶל זִקְנֵי הָעִיר הַשָּׁעְרָה׃}
}
{
	\eDtn{
		then the young woman’s father and her mother shall take
		and bring out the signs of the young woman’s virginity to
		the city’s elders at the gate.
	}
}
{
}

\Verse{16}
{
	\hDtn{וְאָמַר אֲבִי הַנַּעֲרָ אֶל הַזְּקֵנִים אֶת בִּתִּי נָתַתִּי לָאִישׁ הַזֶּה לְאִשָּׁה וַיִּשְׂנָאֶהָ׃}
}
{
	\eDtn{
		“And the young woman’s father shall say to the elders, ‘I
		gave my daughter to this man for a wife, and he hated her.
	}
}
{
}

\Verse{17}
{
	\hDtn{וְהִנֵּה הוּא שָׂם עֲלִילֹת דְּבָרִים לֵאמֹר לֹא מָצָאתִי לְבִתְּךָ בְּתוּלִים וְאֵלֶּה בְּתוּלֵי בִתִּי וּפָרְשׂוּ הַשִּׂמְלָה לִפְנֵי זִקְנֵי הָעִיר׃}
}
{
	\eDtn{
		And, here, he has asserted words of abuse, saying: “I
		didn’t find signs of virginity for your daughter.” But
		these are the signs of my daughter’s virginity!’ And they
		shall spread out the garment in front of the city’s
		elders.
	}
}
{
%	\footnote{
%		garment. It is either the clothing she wore on the wedding night or the
%		cloth or sheet beneath the couple that night, on which a bloodstain
%		would be evidence of her virginity.
%	}
}

\Verse{18}
{
	\hDtn{וְלָקְחוּ זִקְנֵי הָעִיר הַהִוא אֶת הָאִישׁ וְיִסְּרוּ אֹתוֹ׃}
}
{
	\eDtn{
		“And that city’s elders shall take the man and discipline
		him.
	}
}
{
%	\footnote{
%		discipline him. It is unclear if this means a verbal criticism, a
%		flogging, or some other form of chastisement.
%	}
}

\Verse{19}
{
	\hDtn{וְעָנְשׁוּ אֹתוֹ מֵאָה כֶסֶף וְנָתְנוּ לַאֲבִי הַנַּעֲרָה כִּי הוֹצִיא שֵׁם רָע עַל בְּתוּלַת יִשְׂרָאֵל וְלוֹ תִהְיֶה לְאִשָּׁה לֹא יוּכַל לְשַׁלְּחָהּ כּל יָמָיו׃}
}
{
	\eDtn{
		And they shall fine him a hundred weights of silver and
		give it to the young woman’s father because he brought out
		a bad name on a virgin of Israel. And she shall be his for
		a wife: he shall not be able to let her go, all his days.
	}
}
{
%	\footnote{
%		a bad name on a virgin of Israel. Premarital sex is not forbidden
%		elsewhere in the Torah, and men marry women who are not virgins. Only
%		the high priest is absolutely required to marry a virgin (Lev 21:13–14).
%		The issue in this case must therefore be that the woman was represented
%		to be a virgin at the time she was married, and now her husband claims
%		that she was not. Footnote 22:19. to let her go. Meaning: to divorce
%		her. Footnote 22:19. the young woman’s father. Both the woman herself
%		and her father receive some form of compensation. The woman can never be
%		divorced. The father receives monetary compensation.
%	}
}

\Verse{20}
{
	\hDtn{וְאִם אֱמֶת הָיָה הַדָּבָר הַזֶּה לֹא נִמְצְאוּ בְתוּלִים לַנַּעֲרָ׃}
}
{
	\eDtn{
		“But if this thing was true—signs of virginity for the
		young woman were not found—
	}
}
{
}

\Verse{21}
{
	\hDtn{וְהוֹצִיאוּ אֶת הַנַּעֲרָ אֶל פֶּתַח בֵּית אָבִיהָ וּסְקָלוּהָ אַנְשֵׁי עִירָהּ בָּאֲבָנִים וָמֵתָה כִּי עָשְׂתָה נְבָלָה בְּיִשְׂרָאֵל לִזְנוֹת בֵּית אָבִיהָ וּבִעַרְתָּ הָרָע מִקִּרְבֶּךָ׃}
}
{
	\eDtn{
		then they shall take the young woman out to the entrance
		of her father’s house, and the people of her city shall
		stone her with stones so she dies, because she did a
		foolhardy thing in Israel, to whore at her father’s house.
		So you shall burn away what is bad from among you.
	}
}
{
}

\Verse{22}
{
	\hDtn{כִּי יִמָּצֵא אִישׁ שֹׁכֵב עִם אִשָּׁה בְעֻלַת בַּעַל וּמֵתוּ גַּם שְׁנֵיהֶם הָאִישׁ הַשֹּׁכֵב עִם הָאִשָּׁה וְהָאִשָּׁה וּבִעַרְתָּ הָרָע מִיִּשְׂרָאֵל׃}
}
{
	\eDtn{
		“If a man will be found lying with a woman who is a
		husband’s wife, then the two of them shall die: the man
		who lay with the woman, and the woman. So you shall burn
		away what is bad from Israel.
	}
}
{
}

\Verse{23}
{
	\hDtn{כִּי יִהְיֶה נַעֲרָ בְתוּלָה מְאֹרָשָׂה לְאִישׁ וּמְצָאָהּ אִישׁ בָּעִיר וְשָׁכַב עִמָּהּ׃}
}
{
	\eDtn{
		“If it will be that a virgin young woman will be betrothed
		to a man, and a man will find her in the city and lie with
		her,
	}
}
{
}

\Verse{24}
{
	\hDtn{וְהוֹצֵאתֶם אֶת שְׁנֵיהֶם אֶל שַׁעַר הָעִיר הַהִוא וּסְקַלְתֶּם אֹתָם בָּאֲבָנִים וָמֵתוּ אֶת הַנַּעֲרָ עַל דְּבַר אֲשֶׁר לֹא צָעֲקָה בָעִיר וְאֶת הָאִישׁ עַל דְּבַר אֲשֶׁר עִנָּה אֶת אֵשֶׁת רֵעֵהוּ וּבִעַרְתָּ הָרָע מִקִּרְבֶּךָ׃}
}
{
	\eDtn{
		then you shall take the two of them to that city’s gate
		and stone them with stones so they die: the young woman on
		account of the fact that she did not cry out in the city,
		and the man on account of the fact that he degraded his
		neighbor’s wife. So you shall burn away what is bad from
		among you.
	}
}
{
}

\Verse{25}
{
	\hDtn{וְאִם בַּשָּׂדֶה יִמְצָא הָאִישׁ אֶת הַנַּעֲרָ הַמְאֹרָשָׂה וְהֶחֱזִיק בָּהּ הָאִישׁ וְשָׁכַב עִמָּהּ וּמֵת הָאִישׁ אֲשֶׁר שָׁכַב עִמָּהּ לְבַדּוֹ׃}
}
{
	\eDtn{
		But if the man will find the betrothed young woman in the
		field, and the man will take hold of her and lie with her,
		then only the man who lay with her shall die,
	}
}
{
}

\Verse{26}
{
	\hDtn{וְלַנַּעֲרָ לֹא תַעֲשֶׂה דָבָר אֵין לַנַּעֲרָ חֵטְא מָוֶת כִּי כַּאֲשֶׁר יָקוּם אִישׁ עַל רֵעֵהוּ וּרְצָחוֹ נֶפֶשׁ כֵּן הַדָּבָר הַזֶּה׃}
}
{
	\eDtn{
		but you shall not do a thing to the young woman. The young
		woman does not have a sin deserving death, because, just
		as a man would get up against his neighbor and murder him:
		this case is like that;
	}
}
{
}

\Verse{27}
{
	\hDtn{כִּי בַשָּׂדֶה מְצָאָהּ צָעֲקָה הַנַּעֲרָ הַמְאֹרָשָׂה וְאֵין מוֹשִׁיעַ לָהּ׃}
}
{
	\eDtn{
		because he found her in the field, the betrothed young
		woman cried out, and there was no one to save her.
	}
}
{
}

\Verse{28}
{
	\hDtn{כִּי יִמְצָא אִישׁ נַעֲרָ בְתוּלָה אֲשֶׁר לֹא אֹרָשָׂה וּתְפָשָׂהּ וְשָׁכַב עִמָּהּ וְנִמְצָאוּ׃}
}
{
	\eDtn{
		“If a man will find a virgin young woman who is not
		betrothed, and he’ll grasp her and lie with her, and
		they’ll be found,
	}
}
{
}

\Verse{29}
{
	\hDtn{וְנָתַן הָאִישׁ הַשֹּׁכֵב עִמָּהּ לַאֲבִי הַנַּעֲרָ חֲמִשִּׁים כָּסֶף וְלוֹ תִהְיֶה לְאִשָּׁה תַּחַת אֲשֶׁר עִנָּהּ לֹא יוּכַל שַׁלְּחָהּ כּל יָמָיו׃}
}
{
	\eDtn{
		then the man who lay with her shall give the young woman’s
		father fifty weights of silver, and she shall be his for a
		wife. Because he degraded her, he shall not be able to let
		her go, all his days.
	}
}
{
%	\footnote{
%		the young woman’s father. Both the woman herself and her father receive
%		some form of compensation. The woman can never be divorced. The father
%		receives monetary compensation. Footnote 22:29. she shall be his for a
%		wife. Married to the man who raped her?! Presumably, hopefully, this
%		means that the choice was hers as well as her father’s.
%	}
}
